\section{Réalisation des parcours métier}

\subsection{Sprint 6 : exploitation terrain}

\subsubsection{Alertes et qualification}

Une alerte représente un symptôme opérationnel à traiter, et non une intervention.
Elle possède une sévérité, un statut, une station, un contexte technique et une
chronologie. Elle peut provenir automatiquement d'une perte de communication,
d'un défaut OCPP ou d'une règle de maintenance, mais un utilisateur autorisé peut
également créer une alerte manuelle.

L'interface utilise une disposition en deux panneaux : la liste filtrable reste
visible pendant que le détail, les actions et l'historique de l'alerte sélectionnée
sont consultés. Les filtres de sévérité et d'état ne modifient pas la donnée ; ils
sont convertis en paramètres API. Les compteurs de la barre latérale sont dérivés
des notifications non lues liées à ce contexte.

La figure~\ref{fig:interfaces-alerts} présente l'organisation visuelle de ce
centre opérationnel dans le chapitre suivant.

\subsubsection{Interventions, maintenance et SLA}

Une intervention matérialise le travail humain déclenché par une alerte ou une
décision d'exploitation. L'opérateur la crée, choisit une priorité et l'assigne à
un technicien de la même organisation. Le technicien accepte le travail, ajoute
ses notes, documente le diagnostic, les actions et les pièces, joint des photos
\emph{avant/après}, puis soumet son rapport. L'opérateur contrôle ensuite le
résultat et clôt l'intervention.

La maintenance répond à un besoin différent. Elle planifie une action préventive
ou corrective sur une période, indépendamment de l'existence immédiate d'une
alerte. \texttt{MaintenancePlanService} gère le plan, tandis que
\texttt{MaintenanceLifecycleService} contrôle les transitions. Le job
\texttt{GenerateMaintenanceOccurrences} matérialise les occurrences récurrentes.

\texttt{NotificationSlaService} analyse chaque minute les échéances approchantes
ou dépassées. Les notifications sont enregistrées en base, diffusées par Reverb
et éventuellement envoyées par email via la queue. L'utilisateur peut marquer
une notification lue globalement ou dans le contexte de la page consultée.

Les documents d'alerte, d'intervention et de maintenance sont générés par le
backend. Le rapport PDF reprend la référence, les acteurs, les dates, les preuves
et l'historique sans dépendre d'une capture de l'écran. Les pièces jointes passent
par des routes autorisées et sont supprimables selon le cycle de vie de la
ressource.

\subsection{Sprint 5 : recharge, transaction et paiement}

\subsubsection{Assistant de recharge client}

Le client peut ouvrir l'assistant depuis \emph{Start session}, depuis une fiche
station ou depuis la carte. Dans les deux derniers cas, la station est déjà
sélectionnée et le parcours commence directement au choix du connecteur. Le
workflow se compose de cinq étapes :

\begin{enumerate}[leftmargin=8mm]
  \item sélectionner une station disponible ;
  \item choisir un connecteur compatible et libre ;
  \item brancher le câble et attendre la confirmation OCPP ;
  \item autoriser le paiement simulé et vérifier le tarif applicable ;
  \item envoyer la demande de démarrage distant et attendre la transaction OCPP.
\end{enumerate}

Le passage de l'étape de branchement n'est pas manuel. Le frontend interroge la
tentative de recharge et reçoit la mise à jour du connecteur ; il avance uniquement
lorsque le backend confirme une preuve OCPP compatible. Les pictogrammes de prises
aident le client à reconnaître CCS2, Type~2, CHAdeMO ou un autre standard.

\subsubsection{Cycle de session}

\texttt{ChargingAttemptService} conserve le contexte entre les étapes et refuse
une station, un connecteur ou un tarif devenus invalides. Après autorisation,
\texttt{OcppCommandService} demande un
\texttt{RemoteStartTransaction}. Le simulateur répond, puis
\texttt{StartTransaction} devient la preuve de démarrage. Les
\texttt{MeterValues} actualisent l'énergie, la puissance et l'estimation. Le
client peut laisser la vue active en arrière-plan et y revenir depuis la page des
sessions.

\begin{figure}[H]
  \centering
  \begin{tikzpicture}[
    node distance=8mm and 7mm,
    workflowstep/.style={draw=ctline,fill=ctpaper,rounded corners=2mm,
      minimum width=2.65cm,minimum height=1.05cm,align=center,
      font=\small\bfseries,text=ctdark},
    flow/.style={-{Latex[length=2mm]},thick,draw=ctgreen}
  ]
    \node[workflowstep] (select) {Station et\\connecteur};
    \node[workflowstep,right=of select] (plug) {Branchement\\confirmé OCPP};
    \node[workflowstep,right=of plug] (auth) {Autorisation\\de paiement};
    \node[workflowstep,below=of auth] (start) {Transaction\\démarrée};
    \node[workflowstep,left=of start] (meter) {Mesures et\\estimation};
    \node[workflowstep,left=of meter] (stop) {Arrêt, capture\\et reçu};
    \draw[flow] (select) -- (plug);
    \draw[flow] (plug) -- (auth);
    \draw[flow] (auth) -- (start);
    \draw[flow] (start) -- (meter);
    \draw[flow] (meter) -- (stop);
  \end{tikzpicture}
  \caption{Parcours vertical d'une recharge}
  \label{fig:implementation-charging-flow}
\end{figure}

L'arrêt peut être demandé par le client, par un opérateur autorisé, par la borne
ou par une condition technique. \texttt{StopTransaction} clôt la transaction et
fige les valeurs finales. Le coût n'est pas calculé dans React :
\texttt{TariffResolver} sélectionne la règle applicable, puis le backend calcule
les composantes énergie, temps, forfait et remise.

\subsubsection{Adaptateur de paiement simulé}

Le domaine dépend de l'interface \path{PaymentGateway}. Deux implémentations sont
disponibles : \path{SimulatedPaymentAdapter}, utile aux tests unitaires, et
\path{WireMockPaymentAdapter}, qui appelle un serveur HTTP externe local. Les
opérations fonctionnelles sont l'autorisation, la capture, la libération et le
débit direct.

WireMock reproduit des succès, refus, délais, erreurs serveur et webhooks. Le
webhook est signé ; \path{PaymentWebhookSignature} vérifie son authenticité et
\path{PaymentProviderEventService} assure l'idempotence. Un même événement
rejoué ne capture donc pas deux fois une session. Le reçu est consultable dans
l'application et exportable en PDF.

Ce choix ne simule pas une banque au niveau réglementaire. Il valide la frontière
technique qui permettra d'ajouter ultérieurement ClicToPay, e-DINAR/D17,
GPGCheckout, Konnect ou Flouci. L'adaptateur réel devra ajouter la conformité du
prestataire, le rapprochement, la gestion des litiges et des remboursements.

\subsection{Sprint 7 : tarification et commerce}

\subsubsection{Tarifs et plans de recharge}

Un tarif appartient à une organisation et peut combiner un prix au kWh, un prix
au temps, un coût fixe et des règles horaires. Les affectations définissent son
périmètre : organisation, station ou connecteur. Lorsqu'une session commence,
\texttt{TariffResolver} applique l'affectation la plus spécifique et conserve une
référence au tarif retenu afin qu'une modification future ne change pas
rétroactivement le calcul.

Les plans clients sont des avantages proposés par une organisation : remise sur
la recharge, éventuel abonnement périodique et conditions d'accès. Un client peut
cumuler des abonnements auprès de plusieurs réseaux, car son utilisation d'une
borne ne le transforme pas automatiquement en client exclusif de son propriétaire.
La simple recharge crée une relation transactionnelle ; l'abonnement demeure une
décision explicite.

\subsubsection{Abonnement SaaS de l'organisation}

Le contrat commercial de l'organisation est volontairement séparé des paiements
des conducteurs. Il détermine la période d'essai, le nombre maximal d'employés et
de stations, les fonctions accessibles, la date de renouvellement et les factures
adressées à l'organisation. \texttt{OrganizationEntitlementService} contrôle les
limites avant la création d'un actif, tandis que
\texttt{OrganizationSubscriptionLifecycleService} traite essai, activation,
suspension, expiration et restauration.

L'interface de suivi de l'abonnement est reproduite dans la
figure~\ref{fig:interfaces-organization-billing}.

Le super administrateur gère le catalogue SaaS et le portefeuille global.
L'administrateur consulte son offre, sa consommation de capacité et ses factures,
puis émet une demande de changement. Cette demande évite qu'un administrateur
modifie directement un contrat financier. L'encaissement SaaS reste simulé ou
manuel dans le MVP.

\subsection{Sprint 8 : pilotage, reporting et gouvernance}

\subsubsection{Tableaux de bord par rôle}

\texttt{DashboardService} compose une réponse différente selon le rôle plutôt que
de renvoyer le même tableau de bord avec des libellés modifiés. Le super
administrateur supervise les organisations et les intégrations. L'administrateur
observe les employés, clients, revenus, stations et performances de son
organisation. L'opérateur suit disponibilité, incidents et sessions. Le
technicien voit sa charge terrain et ses SLA. Le client consulte ses recharges,
paiements et abonnements.

La figure~\ref{fig:interfaces-admin-dashboard} montre la composition retenue
pour le tableau de bord de l'administrateur.

Les indicateurs utilisent les mêmes règles et périmètres que les écrans de détail.
Par exemple, la disponibilité du tableau de bord passe par
\texttt{AvailabilityMetrics} et ne moyenne pas des valeurs fictives du frontend.
Les sélecteurs de période sont transformés en bornes temporelles côté API.

\subsubsection{Rapports et échanges internes}

Les espaces d'analyse sont adaptés au périmètre du rôle :
\texttt{/reporting/platform}, \texttt{/organization},
\texttt{/operations} et \texttt{/field}. Les exports PDF, JSON et CSV réutilisent
les mêmes filtres et autorisations. Les rapports internes sont des objets
distincts des exports : un employé peut composer un rapport, joindre des fichiers,
choisir un destinataire de son organisation, l'envoyer, puis suivre lecture et
archivage.

La figure~\ref{fig:interfaces-admin-reports} illustre cet espace d'analyse et
d'échange.

\subsubsection{Administration globale}

Le super administrateur dispose d'un espace cohérent avec la navigation du
produit, mais orienté opérations de plateforme : organisations, utilisateurs,
rôles, permissions, demandes de démonstration, catalogue commercial, audit,
intégrations et paramètres. Les secrets OAuth, OCPP, email ou paiement ne sont
jamais renvoyés par les endpoints d'intégration. L'écran expose uniquement l'état
de configuration et la date du dernier contrôle.

\texttt{PlatformSettingService} limite les paramètres modifiables à une liste
validée. Les valeurs sensibles restent dans l'environnement. Toute modification
globale est attribuée à son auteur dans le journal d'audit. Les paramètres
personnels, tels que fuseau horaire, préférences de notification et rayon de
recherche, restent séparés des contrôles globaux.

\subsection{Sprint 9 : exports et évolutions différées}

\texttt{ReportExportService} centralise les formats PDF, CSV et JSON. OpenSpout est
disponible pour les formats tabulaires plus riches, tandis que DomPDF produit les
documents visuels. Les reçus, factures et rapports opérationnels sont générés près
de leur ressource afin que l'utilisateur les retrouve dans le workflow concerné,
et non uniquement dans une page de rapports générique.

La gestion distante du firmware n'est pas simulée comme si elle était achevée.
Elle nécessite une matrice de compatibilité matériel/logiciel, la signature des
paquets, un hébergement sécurisé, le suivi de progression, une stratégie de
retour arrière et une validation sur borne. Elle reste donc une évolution du
Product Backlog non engagée dans le MVP, conformément au risque de 13 points
associé à US-38.
